\section{Mass-Matrix Metric Foundation for Reduced Coordinates}
\label{app:mass_matrix}

This appendix summarizes the metric foundation used in the reduced-coordinate construction. The key object is the $p$-FEM mass matrix, which is the discrete representation of the continuous profile inner product.

Let reconstructed vertical profiles be expanded in basis functions:
\begin{equation}
f(z) = \sum_{i=1}^{N} a_i \phi_i(z),
\qquad
g(z) = \sum_{j=1}^{N} b_j \phi_j(z).
\end{equation}
Define the mass matrix entries by
\begin{equation}
M_{ij} = \int_{z_{\min}}^{z_{\max}} \phi_i(z)\phi_j(z)\,dz.
\end{equation}
Then the coefficient-space bilinear form satisfies
\begin{equation}
\vect{a}^\top \mat{M} \vect{b} = \int_{z_{\min}}^{z_{\max}} f(z)g(z)\,dz,
\end{equation}
so the induced inner product is
\begin{equation}
\langle u,v\rangle_M = \vect{u}^\top \mat{M} \vect{v}.
\end{equation}

Accordingly, norms and orthogonality are measured by
\begin{equation}
\|\vect{u}\|_M^2 = \vect{u}^\top \mat{M} \vect{u},
\end{equation}
\begin{equation}
\vect{u} \perp_M \vect{v} \iff \vect{u}^\top \mat{M} \vect{v} = 0.
\end{equation}
This metric interpretation is central: the mass matrix defines lengths, angles, and overlap in the reduced profile space.

\subsection{Why this matters for artifact defense}

In Euclidean PCA/SVD, the dot product depends directly on discretization geometry (level spacing and clustering), so dominant modes can shift when instrument placement changes. In contrast, the $\mat{M}$-weighted Hilbert-space inner product approximates the continuous overlap integral of profiles over height. As a result, the extracted coordinates emphasize physical profile structure rather than sensor-layout artifacts.

This is the rationale for performing decomposition in the $\mat{M}$-weighted metric:
\begin{equation}
\vect{v}_i^\top \mat{M} \vect{v}_j = \delta_{ij}.
\end{equation}
The resulting basis vectors are orthonormal under physically meaningful overlap, which strengthens cross-campaign comparability and supports mesh-invariant interpretation of reduced coordinates.

\subsection{Connection to evolution equations}

After Galerkin projection, time-dependent reduced models naturally take the form
\begin{equation}
\mat{M}\dot{\vect{a}} = \mat{K}\vect{a} + \vect{F},
\end{equation}
where $\vect{a}$ is the reduced-state coefficient vector, $\mat{K}$ contains projected operators, and $\vect{F}$ is forcing. In this form, $\mat{M}$ is not an auxiliary matrix; it is the geometry-carrying metric operator of the reduced system.

\subsection{Practical implication for SBL transition analysis}

Near strongly stable transitions, local scalar diagnostics can saturate and lose sensitivity. Metric-aware coordinates obtained in the $\mat{M}$-weighted space remain tied to full-column structure and therefore retain geometric signal where one-point thresholds become less informative. This is the mathematical basis for treating the reduced manifold coordinates as physically interpretable state variables rather than purely statistical features.